% Generated draft: requires mathematical and source review.
\documentclass[11pt]{article}
\usepackage[margin=25mm]{geometry}
\usepackage{fontspec}
\setmainfont{Latin Modern Roman}
\usepackage{amsmath,amssymb,mathtools,unicode-math}
\setmathfont{Latin Modern Math}
\usepackage{enumitem,needspace,xurl,hyperref}
\setlength{\parindent}{0pt}
\setlength{\parskip}{5pt}
\setlength{\emergencystretch}{3em}
\newcommand{\N}{\mathbb{N}}
\newcommand{\Z}{\mathbb{Z}}
\newcommand{\Q}{\mathbb{Q}}
\newcommand{\R}{\mathbb{R}}
\newcommand{\C}{\mathbb{C}}
\newcommand{\F}{\mathbb{F}}
\newcommand{\A}{\mathbb{A}}
\newcommand{\PP}{\mathbb{P}}
\newcommand{\E}{\mathbb{E}}
\newcommand{\eps}{\varepsilon}
\newcommand{\dd}{\,\mathrm{d}}
\begin{document}
\section*{UnsolvedMath ID: 800135}
\textbf{Source problem code:} catalogue entry
\subsection{Definitions and mathematical statement}
Let \(\zeta(s)\) be the meromorphic continuation of the Riemann zeta function, initially defined by \(\zeta(s)=\sum_{n=1}^{\infty}n^{-s}\) for \(\operatorname{Re}s>1\). A nontrivial zero is a zero \(\rho=\beta+i\gamma\) with \(0<\beta<1\). A zero is simple if \(\zeta'(\rho)\ne0\). Interpret the source's abbreviation SZ as the hypothesis that every nontrivial zero of \(\zeta\) is simple; the supplied excerpt does not explicitly expand this abbreviation. No Riemann hypothesis is imposed.

The question is to obtain quantitative upper bounds for
\[
\frac{1}{|\zeta'(\rho)|}
\]
as \(\rho\) ranges over the nontrivial zeros. Under SZ, a uniform formulation is to determine an explicit positive function \(U\), with quantitatively controlled growth, such that
\[
\frac{1}{|\zeta'(\beta+i\gamma)|}\le U(|\gamma|)
\qquad\text{for every nontrivial zero }\beta+i\gamma.
\]
The source specifies neither a proposed growth rate for \(U\) nor whether the desired bound is conditional on hypotheses beyond SZ. Thus this is a request for quantitative bounds, not a conjecture asserting a particular asymptotic order. Without simplicity, the reciprocal is not finite at a multiple zero; restricting the question to simple zeros would be a distinct qualification.

For motivation, the source states that, assuming SZ, the Mertens sum satisfies
\[
M_0(x)=\sum_{\rho}\frac{x^{\rho}}{\rho\zeta'(\rho)}-2
+\sum_{n=1}^{\infty}\frac{(-1)^{n-1}(2\pi)^{2n}}{(2n)!\,n\zeta(2n+1)x^{2n}}.
\]
Here \(M_0(x)\) is the source's notation for its Mertens sum. The excerpt does not define its endpoint convention, specify the domain of \(x\), or prescribe the limiting interpretation of the sum over zeros. The displayed formula is retained as supplied context, with these qualifications unresolved.

The source also gives, with a reference labelled [Tao],
\[
\frac{1}{\zeta'(\rho)}
=\lim_{s\to\rho}\frac{s-\rho}{\zeta(s)}
=-\frac{2e^{1-B\rho}\rho(\rho-1)\Gamma(1+\rho/2)\pi^{-\rho/2}}
{\displaystyle\prod_{\rho'\ne\rho}\left(1-\frac{\rho}{\rho'}\right)e^{\rho/\rho'}},
\]
where \(\Gamma\) is the gamma function and \(B=-0.0230957\ldots\) is the constant in the Weierstrass product for the completed zeta function
\[
\xi(s)=\tfrac12s(s-1)\pi^{-s/2}\Gamma(s/2)\zeta(s).
\]
The reciprocal identity concerns a simple zero \(\rho\); the excerpt does not specify a convention for the infinite product. It identifies control of the factors \(1-\rho/\rho'\), and hence lower bounds for distances \(|\rho-\rho'|\) between distinct zeros, as a central difficulty. Its statements that the order of the reciprocal is a major open problem and that extremely close zeros cannot yet be ruled out even under SZ are supplied source assessments, not independently verified literature conclusions.
\subsection{Short English statement}
Obtain explicit quantitative upper bounds for the reciprocal of the derivative of the Riemann zeta function at its nontrivial zeros, uniformly as their imaginary parts grow. The source discusses the simplicity hypothesis SZ and the difficulty caused by very close zeros, but specifies no target growth rate. Its Mertens formula and product identity retain unresolved conventions requiring review.
\subsection{Sources}
\begin{itemize}
\item Local UnsolvedMath registry: supplied complete source record, permanent ID \(800135\), exact source code \texttt{catalogue entry}, title ``Comparative Prime Number Theory --- Problem 4.'' The registry labels the problem open; this is a supplied status, not an independently verified mathematical assessment.
\item Hamieh, Kadiri, Martin, and Ng, \emph{Comparative Prime Number Theory Problem List} (2024), Problem (\#4). Supplied source URL: \url{https://arxiv.org/abs/2407.03530}. The supplied provenance records access on 2026-10-16 and extraction by \texttt{arxiv\_latex}; these are source-record provenance claims, not a claim of access or verification for this draft.
\item The problem text is attributed to Nathan Ng, University of Lethbridge. The registry's proposer field names Hamieh, Kadiri, Martin, and Ng and gives the year 2024.
\item The product identity is attributed in the supplied excerpt to a reference labelled [Tao]. No bibliographic entry or URL for [Tao] is supplied, so this attribution cannot be resolved from the record.
\end{itemize}
\end{document}
